\scanpage{39}
\mspage{38}
\mspage{82}
\margin{n° 382}

variétés sont pures, et si $X'=X$, (resp. $Y'=Y$), l'application $j^k(\mathrm{Id}.,g)$ (resp. $j^k(f,\mathrm{Id}.)$) est un morphisme pour les fibrations sur X (resp. sur Y).

7.1.12. Soient X, Y des variétés de classe $C^r$, $r'\leq r$ et k un entier fini $\leq r'$. Soient X' et Y' les variétés de classe $C^{r'}$ obtenues à partir de X et Y par affaiblissement de structure (n° 4.13.1). Alors $J^k(X',Y')$ est obtenue par affaiblissement de structure de $C^{r-k}$ à $C^{r'-k}$ à partir de $J^k(X,Y)$.

7.2 - \uline{Jets de sections.}

7.2.1. Soient X, Y des variétés de classe $C^r$, $\varpi:Y\longrightarrow X$ une submersion et $k\leq r$. On note $P_x^k(X,\varpi)$ l'ensemble des éléments f de $J_x^k(X,Y)$ tels que $j_y^k(\varpi)\circ f$ soit le jet de l'application identique de X, y étant le but de f. Le jet f a cette propriété si et seulement si c'est le jet d'ordre k d'un germe en x d'application f de X dans Y tel que $\varpi\circ f$ soit le germe de l'application identique.

7.2.2. On suppose dorénavant k fini. L'ensemble $P_x^k(X,\varpi)$ est une sous-variété fermée de $J_x^k(X,Y)$, et la réunion $P^k(X,\varpi)$ des $P_x^k(X,\varpi)$ $(x\in X)$ est une sous-variété fermée de $J^k(X,Y)$. Si X et Y sont des variétés pures, les fibrations de $J^k(X,Y)$ sur $X\times Y$, X et Y (cf. n° 7.1.6) induisent des fibrations de $P^k(X,\varpi)$, de bases respectives $\varpi(X)\times Y$, $\varpi(X)$ et Y. Si $\varpi$ est la projection d'une fibration $\xi$, on écrira aussi $P^k(X,\xi)$ pour $P^k(X,\varpi)$. Si $\varpi$ est la projection d'un fibré vectoriel E, alors $P^k(X,\varpi)$ est un fibré vectoriel sur X, et un fibré en modules sur le fibré en algèbres $P^k(X)$ ; on le notera aussi $P^k(X,E)$. Si de plus E est le fibré vectoriel trivial $F_X$ de fibre F sur X, alors
\[
P^k(X,\varpi)\simeq P^k(X)\otimes F_X\simeq J^k(X,F).
\]

7.2.3. Soit encore k' un entier $\leq k$. L'application $J^k(X,Y)\longrightarrow J^{k'}(X,Y)$ (cf. n° 7.1.2) induit une application $P^k(X,\varpi)\longrightarrow P^{k'}(X,\varpi)$ qui est un morphisme pour les fibrations sur X. Si Y = E est un fibré vectoriel sur X, le noyau du morphisme $P_x^k(X,E)\longrightarrow P_x^{k'}(X,E)$ $(x\in X)$ s'identifie à l'espace vectoriel des jets
